\documentclass[10pt,a4paper]{amsart}
\usepackage[margin=19mm]{geometry}
\usepackage{amsmath,amssymb,mathtools}
\usepackage[scr=rsfs,cal=euler]{mathalfa}
\usepackage{xfrac}
\usepackage[unicode,hidelinks,bookmarks=false]{hyperref}
\newcommand{\ie}{i.e.\ }
\newcommand{\cf}{cf.\ }
\newcommand{\resp}{respectively\ }
\newcommand{\Subsection}{\subsection}
\providecommand{\nobreakdash}{\nobreak\hbox{-}}
\setlength{\emergencystretch}{2em}
\multlinegap0pt
\usepackage{amssymb}
\usepackage{mathtools}
\usepackage{enumerate}
\usepackage{tcolorbox}
\newtheorem{lem}{Lemma}
\newtheorem{thm}{Theorem}
\newcommand{\Aut}{\operatorname{Aut}}
\newcommand{\End}{\operatorname{End}}
\newcommand{\Inn}{\operatorname{Inn}}
\newcommand\ZZ{\mathbb{Z}}
\title{Endo-Twisted Conjugacy and Outer Fixed Points in Solvable Baumslag--Solitar Groups}
\author{Mallika Roy}
\date{Source: arXiv:2602.11031v1 (CC BY 4.0)}
\begin{document}
\maketitle

\noindent\emph{Source and licence.} Endo-Twisted Conjugacy and Outer Fixed Points in Solvable Baumslag--Solitar Groups. Version arXiv:2602.11031v1 (CC BY 4.0), distributed by arXiv under the Creative Commons Attribution 4.0 International licence, http://creativecommons.org/licenses/by/4.0/, as recorded in the arXiv licence metadata for this exact version and shown on the abstract page https://arxiv.org/abs/2602.11031v1. Full licence notice: \texttt{licenses/arxiv-2602.11031-CC-BY-4.0.txt}.\par\smallskip

\noindent\textit{Editorial scope.} Counted unit: the article's thm TCP (source0.tex lines 450--452). The article's complete original proof of it is delivered with it (source0.tex lines 455--584, one \texttt{proof} environment of 130 lines). No mathematics is written, repaired or re-derived: the statement and the proof are the article's own lines, in the article's own notation, and no retained line is substituted or re-flowed.  Retained uncounted as the source-local context the proof needs: lines 250--260; lines 312--325; lines 405--407; lines 409--411; lines 420--428.  Nothing was omitted from any retained span, so the package carries no omission record.  No retained line contains a citation, so the wrapper carries no bibliography and no bibliography entry.  No retained line contains a figure environment, an image-inclusion command or drawing code, so no image is delivered and the package declares no asset dependency.  Editorial formatting only, in addition: an AMS \texttt{amsart} preamble that restates the article's own declarations and macros used by the retained lines, namely the environments \texttt{lem}, \texttt{thm} on separate counters, together with the standard packages those lines need.  The retained environments print in this document as Lemma 1, Theorem 1 in order of appearance, whereas the article prints the same items with the numbering of its own full document.  The complete native source of the article as distributed in the arXiv e-print for this exact version is bundled unchanged as \texttt{sources/194454/source0.tex}.
\begin{lem}\label{lem: calculs}
For any $\alpha =k/n^p,\, \beta=\ell/n^q\in \ZZ[1/n]$ ($k,\ell, p,q\in \ZZ$, $p,q\geq 0$) and $c, r\in \ZZ$, we have 
\begin{itemize}
\item[(i)] $t^c a^{\alpha} =a^{\alpha n^c} t^c$ (equivalently, $a^{\alpha} t^c= t^c a^{\alpha n^{-c}}$);
\item[(ii)] $a^{\alpha}a^{\beta}=a^{\alpha+\beta}$;
\item[(iii)] $(a^{\alpha})^r =a^{r\alpha}$;
\item[(iv)] $t^{-c}a^{\alpha}t^c=a^{\alpha/n^c}$;
\item[(v)] $\big( a^{\alpha}t^c \big)^{-1}=a^{-\alpha n^{-c}} t^{-c}$; 
\item[(vi)] $\big( a^{\alpha}t^c \big)^r = a^{\frac{n^{rc}-1}{n^c-1}\alpha} t^{rc}$.
\end{itemize}
\end{lem}

\begin{lem}\label{lem: Composition}
The composition rules of type-I and type-II endomorphisms are the following:
%
\begin{itemize}
    \item[(i)] $\psi^{I}_{\alpha_1, \beta_1, 1} \circ \psi^{I}_{\alpha_2, \beta_2, 1} =\psi^{I}_{\alpha_1\alpha_2, \beta_1\alpha_2+\beta_2, 1}$;
    
    \item[(ii)] $\psi^{II}_{0,\beta_1, c_1} \circ \psi^{II}_{0,\beta_2, c_2} =\psi^{II}_{0, \mu\beta_2, c_1c_2}$, where $\mu = \frac{n^{c_1c_2}-1}{n^{c_2}-1}$;
    
    \item[(iii)] $\psi^{II}_{0,\beta_1, c_1} \circ \psi^{I}_{\alpha_2, \beta_2, 1} = \psi^{II}_{0, \beta_1\alpha_2 + \mu_1\beta_2, c_1}$, where $\mu_1 = \frac{n^{c_1} - 1}{n-1}$;
    
    \item[(iv)] $\psi^{I}_{\alpha_1, \beta_1, 1} \circ \psi^{II}_{0,\beta_2, c_2} = \psi^{II}_{0,\beta_2, c_2}$.
\end{itemize}
%    
\end{lem}

 \begin{align}\label{eq: action I}
\big( a^{\nu}t^d \big) \psi^I_{\alpha,\beta, 1} = \big( t^{-r}a^{m}t^{r+d}\big) \psi^I_{\alpha,\beta, 1} &= \big( a^{\beta}t\big)^{-r} a^{m\alpha} \big(a^{\beta}t\big)^{r+d} \notag \\ &= a^{\frac{n^{-r}-1}{n-1}\beta}t^{-r} a^{m\alpha} a^{\frac{n^{r+d}-1}{n-1}\beta}t^{r+d} \notag \\ &= a^{\frac{n^{-r}-1}{n-1}\beta} a^{m n^{-r}\alpha} a^{n^{-r}\frac{n^{r+d}-1}{n-1}\beta}t^{d} \\ &= a^{\frac{n^{-r}-1}{n-1}\beta+ m n^{-r}\alpha + \frac{n^{d}-n^{-r}}{n-1}\beta}t^{d} \notag \\ &= a^{\nu \alpha + \frac{n^{d}-1}{n-1}\beta}t^{d}. \notag
 \end{align}

  \begin{align}\label{eq: action II}
\big( a^{\nu}t^d \big) \psi^{II}_{0,\beta, c} = \big( t^{-r}a^{m}t^{r+d}\big) \psi^{II}_{0,\beta, c} &= \big( a^{\beta}t^c\big)^{-r} 1 \big( a^{\beta}t^c\big)^{r+d} \notag \\ &= \big( a^{\beta}t^c\big)^d \\ &= a^{\frac{n^{cd}-1}{n^c-1}\beta}t^{cd}. \notag
 \end{align}

\begin{lem} \label{lem: twisted endos}
Let $G$ be a group, $u,v,w \in G$, $\varphi \in \Aut(G)$, $\psi \in \End(G)$ and $\gamma_w \in \Inn(G)$. Then, the following are equivalent:
 \begin{itemize}
\item[(a)] $u \sim_\psi v$,
\item[(b)] $u \varphi \sim_{\varphi^{-1}\psi\varphi} v \varphi$,
\item[(c)] $wu \sim_{\psi\gamma_{w^{-1}}} wv$, 
\item[(d)] $uw \sim_{\gamma_{w^{-1}}\psi} vw$. \hfill \qed
 \end{itemize}
\end{lem}

\begin{thm}\label{thm: main TCP}
The endo-twisted conjugacy problem is solvable in $BS(1,n)$.
\end{thm}

\begin{proof}
    Let us suppose that we are given a fixed endomorphism $\psi \in \End (BS(1,n))$ and two elements $u = a^{\nu_1}t^{d_1}, v = a^{\nu_1}t^{d_1} \in BS(1, n)$, where $\nu_i =m_i/n^{r_i}\in \ZZ[1/n]$ with $m_i, r_i, d_i\in \ZZ$, $r_i\geq 0$, $i=1,2$. We have to decide whether $a^{\nu_1}t^{d_1}$ and $a^{\nu_2}t^{d_2}$ are $\psi$-twisted conjugated to each other, i.e., whether $u\sim_\psi v$. As there are two types of endomorphisms, we need to consider the following two cases, namely, Case-1: $\psi = \psi^{II}_{0, \beta, c}$, and Case-2: $\psi = \psi^{I}_{\alpha, \beta, 1}$.

    \noindent
    \textit{Proof of Case-1:} Let us take $w=t^{d_1}, \text{ and } w'=t^{-r_1}$. $\psi^{II}_{0, \beta, c}\in \End(BS(1,n))$ is given where $\beta = \ell/n^q$, where $\ell, q, c \in \ZZ$ and $q \geq 0$. We write $\psi^{II}_{0, \beta, c}$ simply as by $\psi^{II}$. Applying Lemma~\ref{lem: twisted endos} (a) $\Leftrightarrow$ (d), 
    %$u \sim_{\psi} v \Leftrightarrow uw^{-1} \sim_{\gamma_w\psi} vw^{-1} \Leftrightarrow a^{\nu_1} \sim_{\gamma_w\psi} a^{\nu_2} t^{d_2 - d_1} \Leftrightarrow t^{-r_1} a^{m_1} t^{r_1}  \sim_{\gamma_w\psi} t^{-r_2} a^{m_2} t^{r_2}t^{d_2 - d_1}$
    %
    \begin{equation}\label{equ: change 1}
        u \sim_{\psi} v \Leftrightarrow uw^{-1} \sim_{\gamma_w\psi} vw^{-1} \Leftrightarrow a^{\nu_1} \sim_{\gamma_w\psi} a^{\nu_2} t^{d_2 - d_1} \Leftrightarrow t^{-r_1} a^{m_1} t^{r_1}  \sim_{\gamma_w\psi} t^{-r_2} a^{m_2} t^{r_2}t^{d_2 - d_1}
    \end{equation}
    %
   Now take $\varphi = \gamma_{w'} \in \Aut(BS(1,n))$. Combining Lemma~\ref{lem: twisted endos} (a) $\Leftrightarrow$ (b) and from~\eqref{equ: change 1} we get the following:
   %
   \begin{equation}
       t^{-r_1} a^{m_1} t^{r_1}  \sim_{\gamma_w\psi} t^{-r_2} a^{m_2} t^{r_2}t^{d_2 - d_1} \Leftrightarrow a^{m_1} \sim_{\Psi^{II}} V,
   \end{equation}
%
where $\Psi^{II} = (\gamma_{w'})^{-1}\gamma_w\psi (\gamma_{w'})$ and $V = (t^{-r_2} a^{m_2} t^{r_2}t^{d_2 - d_1})\varphi = t^{-(r_2 - r_1)}a^{m_2} t^{r_2 - r_1} t^{d_2 - d_1}$. Further, take $r = r_2 - r_1$ and $d = d_2 - d_1$. First, we observe from~\ref{lem: Composition} that any kind of composition of a type-I and a type-II endomorphism, the resulting endomorphism is of type-II.
Thus, adjusting the twisting endomorphism appropriately, we are reduced to checking the endo-twisted conjugation for inputs $a^{m_1}$ and $t^{-r} a^{m_2} t^r t^d$ with $d, m_1, m_2, r \in \ZZ$, $r \geq 0$. Taking the possible $\psi^{II}$-twisted conjugator as $g = a^{\gamma}t^p$, where $\gamma = y/n^x \in \ZZ[1/n]$, $p \in \ZZ$ and using~\eqref{eq: action II} and Lemma~\ref{lem: calculs}(iv) we have the following:
%
\begin{align*}
    (a^{\gamma}t^p \psi^{II})^{-1} \, a^{m_1} \, a^{\gamma}t^p 
    &= (a^{\frac{n^{pc} -1}{n^c -1}\beta}t^{pc})^{-1} \, a^{m_1} \, a^{\gamma}t^p \\
    &= a^{\frac{n^{-pc} -1}{n^c -1}\beta}t^{-pc} \, a^{m_1} \, a^{\gamma}t^p \\
    &= a^{\frac{n^{-pc} -1}{n^c -1}\beta} \, a^{m_1n^{-pc}} \, a^{\gamma n^{-pc}} \, t^{-pc+p}\\
    &= a^{(\frac{1 - n^{pc}}{n^c -1}\beta + m_1 + \gamma)n^{-pc}} \, t^{-pc+p}.
\end{align*}
%
Altogether we have $u \sim_{\psi^{II}} v$ if and only if 
%$a^{m_2/n^r} t^d = a^{(\frac{1 - n^{pc}}{n^c -1}\beta + m_1 + \gamma)n^{-pc}} \, t^{-pc+p}$. 
%
\begin{equation}\label{main endo II-1}
  a^{m_2/n^r} t^d = a^{(\frac{1 - n^{pc}}{n^c -1}\beta + m_1 + \gamma)n^{-pc}} \, t^{-pc+p}.  
\end{equation}
%
Hence, equating both sides of~\eqref{main endo II-1}, firstly, we need $d = p-pc$ and also,
%
\begin{align*}
    \frac{m_2n^{pc}}{n^r} & = \frac{1 - n^{pc}}{n^c -1}\beta + m_1 + \gamma\\
    \Rightarrow m_2n^{pc-r}(n^c - 1) & = (1 - n^{pc})\beta + m_1(n^c - 1) + \frac{y}{n^x}(n^c - 1)\\
    \Rightarrow m_2(n^c - 1)n^{pc+x} & = (1 - n^{pc})\beta n^r n^x + m_1(n^c - 1) n^r n^x + (n^c - 1) n^r y.
\end{align*}
%
Let us apply the change of variable $z = pc + x\in \ZZ$, where $c$ is given; then the last equation of the above computation becomes
%
\begin{align*}
m_2(n^c - 1)n^{ z} & = (1 - n^{z})\beta n^r n^x + m_1(n^c - 1) n^r n^x + (n^c - 1) n^r y \\
\Rightarrow m_2(n^c - 1)n^{ z} & = \beta n^r n^x - \beta n^r n^{ z} + m_1(n^c - 1) n^r n^x + (n^c - 1) n^r y\\
\Rightarrow An^x + By + Cn^{ z} & =0.
\end{align*}
%
From the inputs, namely, $\ell, q, m_1, m_2, r, d, n \in \ZZ$, $q, r \geq 0$, $n \neq \pm 1$, we compute three integers
\[
A = \beta n^r + m_1n^r(n^c-1), \, \, B= n^r(n^c -1), \, \, C = m_2(1-n^c) - \beta n^r.
\]

Summarizing $A,B,C\in \ZZ$; and $u = a^{\nu_1} t^{d_1}$ and $a^{\nu_2} t^{d_2}$ are $\psi^{II}_{0, \beta, c}$-twisted conjugated to each other if and only if the equation 
 \begin{equation}\label{equ: main endo II-2}
An^x +By + Cn^z = 0
 \end{equation}
has an integral solution, for the unknowns $x,y,z\in \ZZ$, with $x\geq 0$. If $B=\pm 1$ this is always the case; so, assume $B\neq \pm 1$.

If $C=0$ it is easy to check: if further $B=0$ then equation~\eqref{equ: main endo II-2} has no solution unless $A=0$; and if $B\neq 0$ compute the finite set $\mathcal{N}=\{1,n,n^2, \ldots \}\subseteq \ZZ/B\ZZ$ (by computing the successive powers of $n$ until obtaining the first repetition modulo $B$) and check whether the set $A\mathcal{N}$ contains $0 \pmod{B}$.

Now assume $C\neq 0$, write $C=C' n^s$ with $s\geq 0$ and $C'\neq 0$ not being multiple of $n$, and note that all possible solutions to equation~\eqref{equ: main endo II-2} have $z \geq -s$. So, with the change of variable $\ZZ \ni z'=z+s$, \eqref{equ: main endo II-2} is equivalent to 
  \begin{equation}\label{equ: main endo II-3}
 An^x +By + C'n^{z'} =0,
  \end{equation}
 for the unknowns $x,y,z'\in \ZZ$, with $x, z'\geq 0$. Let us distinguish two more cases. 

\textbf{Sub-case 1: $B=0$}. Our equation~\eqref{equ: main endo II-3} becomes $An^x=C'' n^{ z'}$ (where, $C'' = - C'$), which has the desired solution if and only if $A/C''$ is a rational number being a (positive or negative) power of $n$. This can be checked by looking at the prime factorizations of $A$, $C'$, and $n$, and so decidable.

\textbf{Sub-case 2: $B\neq 0$}. Now, equation~\eqref{equ: main endo II-3} is equivalent to 
  \begin{equation}\label{equ: main endo II-4}
 An^x \equiv C''n^{ z'} \pmod{B},
 \end{equation}
 for the unknowns $x,z'\in \ZZ$, with $x,z'\geq 0$. Compute the finite set $\mathcal{N}=\{1,n,n^2, \ldots \}\subseteq \ZZ/B\ZZ$. Clearly, equation~\eqref{equ: main endo II-4} admits a solution if and only if the subsets $A\mathcal{N}$ and $C''\mathcal{N}$ from $\ZZ/B\ZZ$ intersect non-trivially, which is  decidable. This completes the proof of Case-1.

\noindent
\textit{Proof of Case-2:} Let $\psi^{I}_{\alpha, \beta, 1} \in \End(BS(1,n))$, where $\alpha = k/n^j$, $\beta= \ell/ n^q$ with $k, j, \ell, q \in \ZZ$, $j, q \geq 0$, be the given endomorphism. We observe that any type-I endomorphism preserve the $|.|_t$. Hence, $d_1=d_2$ (call it just $d$) is an obvious necessary condition for $u \sim_{\psi^{I}_{\alpha, \beta, 1}} v$. 

Applying Lemma~\ref{lem: twisted endos} (a)$\Leftrightarrow$(d) with $w=t^{-d}$ and (a)$\Leftrightarrow$(b) with $\varphi=\gamma_{t^{-r_1}}$, we have that 
 $$
a^{\nu_1}t^d \sim_{\psi^I_{\alpha, \beta, 1}} a^{\nu_2}t^d \quad \Leftrightarrow \quad a^{\nu_1} \sim_{\gamma_{w^{-1}}\psi^I_{\alpha, \beta,1}} a^{\nu_2} \quad \Leftrightarrow \quad a^{m_1} \sim_{\varphi^{-1}\gamma_{w^{-1}}\psi^I_{\alpha, \beta, 1}\varphi} t^{-(r_2-r_1)}a^{m_2}t^{r_2-r_1}.
 $$
Thus, adjusting the twisting endomorphism appropriately, we are reduced to check twisted conjugation for inputs of the form $u=a^{m_1}$, $v=t^{-r}a^{m_2}t^{r}$, with $m_1, m_2, r (=r_2 - r_1)\in \ZZ$, $r\geq 0$.

Similarly, as in Case-1, writing the possible $\psi^I_{\alpha, \beta, 1}$-twisted conjugator as $g=a^{\gamma}t^p$, $\gamma\in \ZZ[1/n]$, $p\in \ZZ$, and using equation~\eqref{eq: action I}, we have 
 \[
g\psi^I_{\alpha,\beta, 1} =a^{\gamma \alpha + \frac{n^{p}-1}{n-1}\beta}t^{p}; 
\] 
now, using Lemma~\ref{lem: calculs}(iv),  
 \[
\big( g\psi^I_{\alpha,\beta, 1} \big)^{-1} ug=t^{-p} a^{-\gamma \alpha - \frac{n^{p}-1}{n-1}\beta} a^{m_1} a^{\gamma}t^p= a^{n^{-p}\big(-\gamma \alpha -\frac{n^p-1}{n-1}\beta +m_1 +\gamma \big)}.
 \]
Hence, $u=a^{m_1}\sim_{\psi^I_{\alpha, \beta, 1}} t^{-r}a^{m_2}t^r =v$ if and only if 
 \[
\frac{m_2}{n^{r}} n^p =-\gamma \alpha -\frac{n^{p}-1}{n-1}\beta +m_1 +\gamma, 
 \]
for some $\gamma \in \ZZ[1/n]$ and $p\in \ZZ$; rearranging and simplifying, we get
 \[
(n-1)m_2 n^p +(n-1)n^{r}\gamma (\alpha -1) +(n^{p}-1)n^{r}\beta =(n-1)n^{r}m_1.
 \]
Write $\gamma =y/n^x$, $x,y\in \ZZ$, $x\geq 0$, and apply the change of variable $\ZZ \ni z=p+x$; our equation becomes equivalent to the integral equation  
%  \[
% (n-1)m_2 n^p +(n-1)n^{r}\frac{y}{n^x} \frac{k-n^j}{n^j} +(n^{p}-1)n^{r}\frac{\ell}{n^q} =(n-1)n^{r}m_1, 
%  \]
%  \[
% (n-1)m_2 n^{p+x+j+q} +(n-1)y(k-n^j)n^{r+q} +(n^{p}-1)\ell n^{r+j+x} =(n-1)m_1 n^{r+j+q+x}, 
%  \]
%  \[
% \Big( \ell n^{r+j}+(n-1)m_1 n^{r+j+q}\Big) n^x +\Big( (n-1)(n^j-k)n^{r+q} \Big) y =\Big( (n-1)m_2 n^{j+q} +\ell n^{r+j}\Big) n^{z}.
%  \]
 %
\begin{align*}
 (n-1)m_2 n^p +(n-1)n^{r}\frac{y}{n^x} \frac{k-n^j}{n^j} +(n^{p}-1)n^{r}\frac{\ell}{n^q} & =(n-1)n^{r}m_1,\\ 
 (n-1)m_2 n^{p+x+j+q} +(n-1)y(k-n^j)n^{r+q} +(n^{p}-1)\ell n^{r+j+x} &=(n-1)m_1 n^{r+j+q+x},\\
\Big( \ell n^{r+j}+(n-1)m_1 n^{r+j+q}\Big) n^x +\Big( (n-1)(n^j-k)n^{r+q} \Big) y &=\Big( (n-1)m_2 n^{j+q} +\ell n^{r+j}\Big) n^{z}. 
\end{align*}
% 
From the data, namely $k, j, \ell, q, m_1, m_2, r, n\in \ZZ$, $j,q,r\geq 0$, $n\neq \pm 1$, let us compute the three integers 
 \[
A=\ell n^{r+j}+(n-1)m_1 n^{r+j+q},\qquad B=(n-1)(n^j-k)n^{r+q}, \qquad C=(n-1)m_2 n^{j+q}+\ell n^{r+j}.
 \]
Thus, we get  $A,B,C\in \ZZ$; and $u=a^{m_1}$ and $v=t^{-r}a^{m_2}t^r$ are $\psi^I_{\alpha, \beta, 1}$-twisted conjugated to each other if and only if the equation 
 \begin{equation}\label{eq: equation-I 1}
An^x +By=Cn^z
 \end{equation}
has an integral solution, for the unknowns $x,y,z\in \ZZ$, with $x\geq 0$. If $B=\pm 1$ this is always the case; so, assume $B\neq \pm 1$. Then, we execute the same procedure as of the proof of the Case-1 with these new values of $A, B$ and $C$, in particular we study all the sub-cases, e.g., $C=0$ with $B = 0$ or $B \neq 0$, and $C \neq 0$ with $B = 0$ or $B \neq 0$. And, we see that all of these sub-cases are decidable. This completes the proof of Case-2. 
\end{proof}

\end{document}
